\documentclass{article}

\usepackage{arxiv}

\usepackage{booktabs}
\usepackage{array}
\usepackage[utf8]{inputenc} % allow utf-8 input
\usepackage[T1]{fontenc}    % use 8-bit T1 fonts
\usepackage{lmodern}        % https://github.com/rstudio/rticles/issues/343
\usepackage{hyperref}       % hyperlinks
\usepackage{url}            % simple URL typesetting
\usepackage{booktabs}       % professional-quality tables
\usepackage{amsfonts}       % blackboard math symbols
\usepackage{nicefrac}       % compact symbols for 1/2, etc.
\usepackage{microtype}      % microtypography
\usepackage{graphicx}

\title{Distance, Mobility, and Mode Choice : A Spatial Analysis of
School Travel Behaviour in Ireland}

\author{
    Sananthaa Jagadesan Sethamaraikannan
   \\
    Department of Mathematics \& Statistics \\
    University of Galway \\
  Galway,Ireland. \\
  \texttt{\href{mailto:sananthaajagadesan@gmail.com}{\nolinkurl{sananthaajagadesan@gmail.com}}} \\
   \And
    Judith Beltran Lopez
   \\
    Department of Mathematics \& Statistics \\
    University of Galway \\
  Galway,Ireland \\
  \texttt{} \\
  }


% tightlist command for lists without linebreak
\providecommand{\tightlist}{%
  \setlength{\itemsep}{0pt}\setlength{\parskip}{0pt}}



\providecommand{\pandocbounded}[1]{#1}
\usepackage{booktabs}
\usepackage{array}
\providecommand{\pandocbounded}[1]{#1}
\begin{document}
\maketitle


\begin{abstract}
In Ireland, whether a child can walk or cycle to school depends heavily
on where they live, something that has been described as a postcode
lottery. This study investigates whether the location of schools and how
easily they can be reached explains differences in active travel rates
across Ireland. We used data from the 2022 Census, the Department of
Education, and OpenStreetMap to measure school accessibility for over
16,000 small areas across Ireland. Two methods were used to measure
accessibility-kriging, which estimates the density of school places
across space, and r5r, which calculates how many school places can be
reached within a given walking or cycling time using the real road
network. We then used Bayesian Beta regression to model the relationship
between accessibility and the proportion of children walking or cycling
to school.
\end{abstract}

\keywords{
    -sss
  }

\section{Introduction}\label{introduction}

Here goes an introduction text

\subsection{1.1 Background}\label{background}

\label{sec:headings}

You can use directly LaTeX command or Markdown text.

LaTeX command can be used to reference other section. See Section
\ref{sec:headings}. However, you can also use \textbf{bookdown}
extensions mechanism for this.

\subsection{1.2 The Postcode Lottery Argument from the
article}\label{the-postcode-lottery-argument-from-the-article}

\subsection{1.3 Research Objectives}\label{research-objectives}

\section{2 Data}\label{data}

\subsection{2.1 School Data}\label{school-data}

\subsection{2.2 Census Data}\label{census-data}

\subsection{2.3 Study Area}\label{study-area}

\section{3 Methods}\label{methods}

\subsection{3.1 School Place
Density:Kriging}\label{school-place-densitykriging}

\subsection{3.2 Travel Time
Accessibility:r5r}\label{travel-time-accessibilityr5r}

\subsection{3.3 Demographic Variables:ILR
Transformation}\label{demographic-variablesilr-transformation}

\subsection{3.4 Statistical Model:Bayesian Beta
Regression}\label{statistical-modelbayesian-beta-regression}

\subsection{3.5 Normalisation by School:Age
Population}\label{normalisation-by-schoolage-population}

\section{4 Results}\label{results}

\subsection{4.1 Spatial Distribution of
Schools}\label{spatial-distribution-of-schools}

\subsection{4.2 Walking and Cycling Accessibility: travel
time}\label{walking-and-cycling-accessibility-travel-time}

\subsection{4.3 Kriging vs Travel Time
Comparison}\label{kriging-vs-travel-time-comparison}

\subsection{4.4 Model Results --- Urban Areas,rural if
possible}\label{model-results-urban-areasrural-if-possible}

\subsection{4.5 Effect of Walking Accessibility on Active
Travel}\label{effect-of-walking-accessibility-on-active-travel}

\section{5 Discussion}\label{discussion}

\subsection{5.1 Interpretation of
Results}\label{interpretation-of-results}

\subsection{5.2 Limitations}\label{limitations}

\subsection{5.3 Comparison with the
article}\label{comparison-with-the-article}

\section{6 Conclusions}\label{conclusions}

\section{7 References}\label{references}

\bibliographystyle{unsrt}
\bibliography{references.bib}


\end{document}
